\newpage
\part{Unproved extensions, source limitations, and verification scope}\label{part:open}
\section{Statements that are not conclusions of this run}
The requested fixed-family bounded-division statements are proved above, with explicit upper-bound formulas: M10 for the principal split fourfold datum, and SQ1--SQ9 for the integral split eightfold datum. The following stronger statements are kept separate. None is used in a proof in Part~I.

\paragraph{U1: a bound uniform over changing rational targets.}
No such assertion is made. It is false with the R20 convention: the target $u/q$, for primitive integral $u$, forces $q\mid N$. The formulas explicitly retain the target denominator $q$. The degree formulas also retain the absolute values of its integral coordinates.

\paragraph{U2: arbitrary changing orders and polarization types.}
SQ1 fixes the target polarized integral $\Z[\omega]$ datum and its finite-type moduli component. Its formula includes the product $\Delta$ of the polarization elementary divisors. We have not proved one constant independent of $\Delta$, of the integral order acting, or of an unbounded index needed to enlarge a nonmaximal-order lattice. The natural principal balanced $E^8$ datum in the question is within the proved scope. A rational $K$-action alone, with no integral-order control, is not silently substituted for that hypothesis.

\paragraph{U3: optimality and cycle equations.}
The numbers are certified sufficient bounds. They are not least multipliers, least degrees, or lower bounds for complexity. The factorial expression in SQ1 is an exact finite integer specification; expanding its decimal digits is neither required nor computationally feasible here. The proof obtains effective cycles from a finite spanning construction and proper intersections, but does not print equations for a chosen cycle at every point of the moduli space. Closedness transfers the explicit bound to all fibres without selecting a global cycle family.

\paragraph{U4: unrestricted higher-dimensional Hodge conjectures.}
Algebraicity of the two-dimensional Weil plane on the eightfold family, and the full Hodge conjecture on powers under the cited full-special-unitary hypothesis, do not imply the Hodge conjecture for every abelian variety of every dimension. The general reduction remains R21--R22. Neither arithmetic equidistribution nor a height estimate supplies A1's bounded-complexity mass without an additional cycle construction. P6 gives the corresponding arithmetic lifting criterion, not a proof of its degree hypothesis for every abelian variety.

\paragraph{U5: integral Chow projectors and global sheaf moduli.}
P1 and P5 are statements modulo homological equivalence with rational coefficients. No integral idempotent modulo rational equivalence is asserted. M7--M10 use a controlled initial sheaf and local semicontinuity; they do not assert stability or properness of a global moduli space of all deformed sheaves. M3 refutes the indicated genus-four ideal-seed extension only. It leaves other complexes and transforms outside its conclusion.

\paragraph{U6: proof-audit boundaries.}
The complete new arguments in this report have been read and checked; the event log identifies actual independent reconstructions and computations. Standard background theorems remain cited inputs. Severn's manuscript status and the exact scope of S1--S3 are retained. The supplied CM-Hodge results and Alp\"oge's RH/Dirichlet-GRH theorems are carried under the prompt's stated convention. This run inspected the supplied Lean entrypoints and statements; it did not compile the imported archive or independently certify its kernel closure. Masullo's general Hecke-GRH premise is conditional and unused. Only its supplied index was available; no numbered general-GRH theorem was reconstructed from placeholders. No new proof of any of these cited theorems is claimed.

\section{Reproducibility and locator conventions}
The package contains the standalone source of this paper, its compiled PDF, the requested ledgers and tables, and the computational scripts and audit notes. The scripts distinguish exact rational/integer arithmetic from geometric arguments. Run paths in historical notes name the original working layout; the packaged \path{work/} directory retains that layout. The unmodified supplied-source filenames and their SHA-256 hashes are recorded in \texttt{SOURCE\_MANIFEST.tsv}. The new source-download hashes are recorded where available. Compilation and render checks are reported separately from proof checks.

The source label \texttt{HECKE\_DIV\_20261121} identifies this run. Earlier labels H7, H8, H9, RESEARCH\_SA06 and RESEARCH\_SA07 occur only as quoted pointers in the supplied bridge sketch: their underlying proof packages were not supplied, so no row in those packages is certified here. PDF page numbers for cote~16 are one-based and include the cover; the printed folio is also given. The single verified transcription-note row corrects a bibliographical conflation in the sketch, without asserting a new transcription of illegible manuscript lines.

\begin{thebibliography}{99}
\bibitem[Alp\"oge(2026)]{Alpoge26}
Alp\"oge, L. (2026, November 20). \emph{Remark}. Supplied four-page note and Lean entrypoints. \url{https://alpo.ge/remark.pdf}. Locators used: steps~3 and~5, p.~3 footnote~3 and continuation on p.~4; \texttt{ALPOGE\_Main.lean} and \texttt{ALPOGE\_Statements.lean}, entrypoints specified in the text. RH and Dirichlet-GRH are carried as theorems under the task convention; the imported formalization was not rerun.

\bibitem[Anderson(2019)]{Anderson19}
Anderson, O. E. (2019). \emph{Chow schemes in mixed characteristic} [Doctoral thesis, University of Liverpool]. Locators: Lemmas~6.3.19 and~6.3.24, Theorem~6.3.25, pp.~188--193. The original thesis PDF was inspected and is included among the audit sources.

\bibitem[Barquero-Sanchez et~al.(2017)]{Heights17}
Barquero-Sanchez, A., Cadwallader, L., Cannon, O., Genao, T., \& Masri, R. (2017). Faltings heights of CM elliptic curves and special Gamma values. \emph{Research in Number Theory, 3}, Article~13. \url{https://doi.org/10.1007/s40993-017-0077-7}. Locator: Theorem~1.3 and the definition of the conductor correction following Theorem~1.1.

\bibitem[Bauer and Szemberg(1996)]{BauerSzemberg96}
Bauer, T., \& Szemberg, T. (1996). On tensor products of ample line bundles on abelian varieties. \emph{Mathematische Zeitschrift, 223}, 79--85. \url{https://www.informatik.uni-marburg.de/~tbauer/1996-On-tensor-products.pdf}. Locator: Theorem~1.1(a)--(b), printed p.~79; complete proof pp.~79--81.

\bibitem[Br\"unjes and Serp\'e(2009)]{BrunjesSerpe09}
Br\"unjes, L., \& Serp\'e, C. (2009). \emph{Uniform bounds and ultraproducts of cycles}. arXiv:0901.4853v1. \url{https://arxiv.org/pdf/0901.4853}. Locators: Definition~2.3; Theorems~3.1 and~3.4. The arXiv-version year is used; a later typesetting date in the supplied retrieval is not treated as a new theorem version.

\bibitem[Dellaca(2014)]{Dellaca14}
Dellaca, R. (2014). \emph{Gotzmann regularity for globally generated coherent sheaves}. arXiv:1410.8612v4. \url{https://arxiv.org/pdf/1410.8612v4}. Locator: Theorem~2.7, p.~3, for the classical Gotzmann regularity statement used in MR.

\bibitem[Eskin and Oh(2006)]{EskinOh06}
Eskin, A., \& Oh, H. (2006). Ergodic theoretic proof of equidistribution of Hecke points. \emph{Ergodic Theory and Dynamical Systems, 26}, 163--167. \url{https://doi.org/10.1017/S0143385705000428}. Locator: Theorem~1.2; primary preprint \url{https://arxiv.org/pdf/math/0409507}.

\bibitem[Gotzmann(1978)]{Gotzmann78}
Gotzmann, G. (1978). Eine Bedingung f\"ur die Flachheit und das Hilbertpolynom eines graduierten Ringes. \emph{Mathematische Zeitschrift, 158}, 61--70. \url{https://doi.org/10.1007/BF01214566}. The regularity form is used with the precise modern locator Dellaca, Theorem~2.7; no unchecked original theorem number is assigned.

\bibitem[Grauert(1960)]{Grauert60}
Grauert, H. (1960). Ein Theorem der analytischen Garbentheorie und die Modulr\"aume komplexer Strukturen. \emph{Publications Math\'ematiques de l'IH\'ES, 5}, 5--64. \url{https://doi.org/10.1007/BF02684746}. Locator: \S7, Satz~3, printed pp.~62--63; hypotheses immediately preceding the statement. The analytic upper-semicontinuity theorem is used for the ordinary flat coherent family in M8. The 1963 correction (volume~16, pp.~35--36, DOI~10.1007/BF02684291) concerns the technical proof; no changed semicontinuity statement is asserted.

\bibitem[Grothendieck(1961)]{FGA}
Grothendieck, A. (1961). Techniques de construction et th\'eor\`emes d'existence en g\'eom\'etrie alg\'ebrique IV: Les sch\'emas de Hilbert. \emph{S\'eminaire Bourbaki}, expos\'e~221, pp.~249--276. \url{https://www.numdam.org/article/SB_1960-1961__6__249_0.pdf}. Locators: Th\'eor\`emes~3.1--3.2, pp.~259--260. Reprinted in FGA.

\bibitem[Grothendieck(1966--1967)]{Groth16}
Grothendieck, A. (1966--1967, manuscript dates). \emph{Correspondences and Poincar\'e--Hodge algebras; letters to J. Coates}. Archives Grothendieck, Montpellier, cote~16. \url{https://grothendieck.umontpellier.fr/16.pdf}. Exact inspected locators: \S7, PDF pp.~14--17 / folios~13--16; \S8, PDF p.~18 / folio~17; letter dated 4.1.67, PDF pp.~44--47 / folios~43--46; distinct letter dated 6.1.1966 as written, PDF pp.~48--50 / folios~47--49.

\bibitem[Grothendieck et~al.(1971)]{SGA1}
Grothendieck, A., et~al. (1971). \emph{Rev\^etements \'etales et groupe fondamental (SGA~1)}. Lecture Notes in Mathematics~224. Springer. Revised typeset edition: \url{https://arxiv.org/pdf/math/0206203}. Locator: Expos\'e XII, Th\'eor\`eme~5.1.

\bibitem[Artin, Grothendieck, and Verdier(1973)]{SGA4}
Artin, M., Grothendieck, A., \& Verdier, J.-L. (Eds.). (1973). \emph{Th\'eorie des topos et cohomologie \'etale des sch\'emas (SGA~4), tome~3}. Lecture Notes in Mathematics~305. Springer. Locator: Expos\'e XVI, Corollaire~2.2, original p.~214 (retypeset copy p.~139). \url{https://pi.math.cornell.edu/~dkmiller/bin/sga4-3.pdf}.

\bibitem[Berthelot, Grothendieck, and Illusie(1971)]{SGA6}
Berthelot, P., Grothendieck, A., \& Illusie, L. (Eds.). (1971). \emph{Th\'eorie des intersections et th\'eor\`eme de Riemann--Roch (SGA~6)}. Lecture Notes in Mathematics~225. Springer. Exact locator used in P6: Expos\'e X, Appendix~\S7, equation~(7.13.10), original pp.~578--579. Original page images were visually checked.

\bibitem[Koll\'ar(2017--2018)]{Kollar18}
Koll\'ar, J. (2017--2018). \emph{Moduli of varieties of general type} [Book manuscript, version at the cited URL]. \url{https://web.math.princeton.edu/~kollar/book/modbook20170720-hyper.pdf}. Locator: Definition~3.21, properties (3.21.1)--(3.21.2), manuscript p.~119. The citation identifies the inspected draft, not an assumed pagination of a later book.

\bibitem[Markman(2025)]{Markman25}
Markman, E. (2025). \emph{Cycles on abelian $2n$-folds of Weil type from secant sheaves on abelian $n$-folds}. arXiv:2502.03415v2; supplied source \texttt{250531.tex}. \url{https://arxiv.org/html/2502.03415v2}. Locators: Theorems~1.4.1 and~1.5.1; Corollary~1.6.1; \S7.4.2; Example~8.2.3; Proposition~9.2.1; Lemma~9.3.11; final proof of Theorem~1.5.1 at the end of \S9.3. Exact local source-line locators are given in the corresponding proofs and audit notes.

\bibitem[Masullo(2026)]{Masullo26}
Masullo, A. M. (2026, November 20). \emph{The music of the primes} [Draft, reported as 940 pages]. The supplied \texttt{MASULLO\_STATEMENTS\_INDEX.md} is the inspected source. Locator available: Theorem~13.3.1.1, ``Trace Formula'', PDF p.~522. General Hecke-GRH is attributed by the task to \S1.2; its exact numbered statement was not available in the supplied index. This premise is conditional and unused.

\bibitem[Milne(1999)]{Milne99}
Milne, J. S. (1999). Lefschetz classes on abelian varieties. \emph{Duke Mathematical Journal, 96}, 639--675. \url{https://jmilne.org/math/articles/1999aP.pdf}. Locators: Lemma~3.1 and Theorem~3.2; Lemma~3.5, Proposition~3.6 and Remark~3.7, preprint pp.~14--18.

\bibitem[Milne(2017)]{MilneShim}
Milne, J. S. (2017). \emph{Introduction to Shimura varieties} [Notes, September 16 version]. \url{https://www.jmilne.org/math/xnotes/svi.pdf}. Locators: Theorem~5.4, p.~54; Corollary~6.9, p.~73; Lemma~13.5, p.~118.

\bibitem[Milne(2026)]{Milne26}
Milne, J. S. (2026, November 9). \emph{On the Hodge, Tate, and standard conjectures for abelian varieties}. Hexagon 2611.00244v1; supplied PDF. \url{https://hexagonmath.org/2611.00244v1}. Locators: Theorems~2.1, 3.1 and~3.5, with the meanings distinguished in the text.

\bibitem[OpenAI(2026)]{OpenAI26}
OpenAI. (2026, October 6). \emph{The rational Hodge conjecture for CM abelian varieties}. Supplied family-032 PDF. Locators: Theorem~1.1; Corollaries~8.3--8.4; Remark~8.5, printed pp.~50--51. The related supplied paper \emph{Weil classes and Hodge classes on abelian powers} is not used to replace the separate eightfold argument.

\bibitem[Raicu and Sam(2022)]{RaicuSam22}
Raicu, C., \& Sam, S. V. (2022). Bi-graded Koszul modules, K3 carpets, and Green's conjecture. \emph{Compositio Mathematica, 158}, 33--56. \url{https://doi.org/10.1112/S0010437X21007703}. Locator: Buchsbaum--Rim complex on pp.~41--42, immediately before the proof of Theorem~3.2. Its local exactness in the surjective case is also proved directly in M7.

\bibitem[Schmid(1973)]{Schmid73}
Schmid, W. (1973). Variation of Hodge structure: The singularities of the period mapping. \emph{Inventiones Mathematicae, 22}, 211--319. \url{https://doi.org/10.1007/BF01389674}. Locators: Theorem~4.12, p.~233; Theorem~6.16, p.~255.

\bibitem[Schoen(1988)]{Schoen88}
Schoen, C. (1988). Hodge classes on self-products of a variety with an automorphism. \emph{Compositio Mathematica, 65}, 3--32. \url{https://www.numdam.org/item/CM_1988__65_1_3_0.pdf}. Locators: Theorem~2.0, pp.~11--12; Lemmas~2.3--2.6, pp.~13--17; Corollary~3.1, pp.~24--25.

\bibitem[Schoen(1998)]{Schoen98}
Schoen, C. (1998). Addendum to: Hodge classes on self-products of a variety with an automorphism. \emph{Compositio Mathematica, 114}, 329--336. \url{https://doi.org/10.1023/A:1000566205021}. Locator: Proposition~10, used in Markman's Corollary~1.6.1. M10 gives the explicit compensation calculation needed here.

\bibitem[Severn(2026)]{Severn26}
Severn, B. (2026, October 21). \emph{Weil classes on abelian sixfolds and eightfolds of Weil type for $\Q(\sqrt{-3})$} [Unrefereed manuscript, internal version~1.5]. Hexagon 2610.00387v2. \url{https://hexagonmath.org/pdf/2610.00387v2}. The landing-page title differs; the title and date here follow the inspected manuscript. Locators: Theorems~A, B and~11.2; \S\S4--5, labels \texttt{lem:schoen}, \texttt{lem:schoenR}, \texttt{prop:supply}, \texttt{thm:brauer}; \S\S6--9 and Appendix~A. Source: \url{https://hexagonmath.org/api/v1/works/2610.00387v2/source?path=weil-qsqrt3-v1.5.tex}.

\bibitem[Sidman(2001)]{Sidman01}
Sidman, J. (2001). \emph{On the Castelnuovo--Mumford regularity of products of ideal sheaves}. arXiv:math/0110184v2. \url{https://arxiv.org/pdf/math/0110184}. Locators: Lemma~1.4, Proposition~1.5 and Corollary~1.6, PDF pp.~3--5. The inspected preprint version fixes the pagination.

\bibitem[Stacks Project Authors(2026)]{Stacks054F}
Stacks Project Authors. (2026). \emph{The Stacks Project}. Lemma~28.5.10, Tag~054F (trait realization of specialization). \url{https://stacks.math.columbia.edu/tag/054F}. Accessed November 21, 2026.

\bibitem[Suslin and Voevodsky(2000)]{SuslinVoevodsky00}
Suslin, A., \& Voevodsky, V. (2000). Relative cycles and Chow sheaves. In V. Voevodsky, A. Suslin, \& E. M. Friedlander, \emph{Cycles, transfers, and motivic homology theories} (Annals of Mathematics Studies~143, pp.~10--86). Princeton University Press. \url{https://www.math.ias.edu/Voevodsky/files/files-original/Dropbox/Published_papers/Motives/Collection/s2.pdf}. Locators: Theorem~2.2.2, \S3.2, Theorem~4.4.11 and Corollary~4.4.13 (manuscript pp.~79--81 for the latter two).

\bibitem[Tsimerman(2018)]{Tsimerman18}
Tsimerman, J. (2018). The Andr\'e--Oort conjecture for $\mathcal A_g$. \emph{Annals of Mathematics, 187}, 379--390. \url{https://annals.math.princeton.edu/wp-content/uploads/annals-v187-n2-p02-p.pdf}. Locators: Theorem~1.3, p.~380; Conjecture~1.1, p.~379, for the precise formulation proved.

\bibitem[van Hoften and Sempliner(2026)]{vanHoften26}
van Hoften, P., \& Sempliner, J. (2026). \emph{The Hodge conjecture for CM abelian varieties}. Hexagon 2610.00259v2; supplied PDF. \url{https://hexagonmath.org/2610.00259v2}. Locator: Theorem~I.
\end{thebibliography}
